\documentclass[10pt,aspectratio=169]{beamer}
\usetheme{Madrid}
\usecolortheme{whale}

\usepackage[utf8]{inputenc}
\usepackage[french]{babel}
\usepackage{graphicx}
\usepackage{booktabs}
\usepackage{xcolor}
\usepackage{amsmath}

\title[Modèle Linéaire RLS-ARX]{Modélisation Adaptative de la Disponibilité en Calcul Volontaire}
\subtitle{Le Modèle Linéaire RLS-ARX : Explication Conceptuelle, Frugalité et Résilience aux Délestages}
\author[Mémoire de Master 2]{Présenté par : \textbf{Neussi} \\[0.5em] \small Encadrant : \textbf{Dr. Adamou Hamza}}
\institute[UY1]{Université de Yaoundé I \\ Département d'Informatique \\ Spécialité : Systèmes et Réseaux}
\date[Juin 2026]{Soutenance de Travaux de Recherche}

\begin{document}

% Slide de Titre
\begin{frame}
  \titlepage
\end{frame}

% Slide Plan
\begin{frame}{Plan de la Présentation}
  \begin{columns}[T]
    \begin{column}{0.48\textwidth}
      \begin{block}{Piliers Théoriques}
        \begin{itemize}
          \item \textbf{1. Limites des approches classiques}
          \item \textbf{2. Concept du modèle ARX}
          \item \textbf{3. Algorithme RLS en ligne}
        \end{itemize}
      \end{block}
      
      \begin{block}{Adaptation locale}
        \begin{itemize}
          \item \textbf{4. Équivalence temporelle de 72h}
          \item \textbf{5. Préservation des connaissances}
        \end{itemize}
      \end{block}
    \end{column}
    
    \begin{column}{0.48\textwidth}
      \begin{block}{Validation et Synergie}
        \begin{itemize}
          \item \textbf{6. Analyse comparative} (OLS vs RLS)
          \item \textbf{7. Architecture prédictive hybride}
          \item \textbf{8. Frugalité matérielle} (TinyML)
          \item \textbf{9. Conclusion et perspectives}
        \end{itemize}
      \end{block}
    \end{column}
  \end{columns}
\end{frame}

% Slide 1: Limites des Approches Classiques
\begin{frame}{1. Pourquoi abandonner les modèles classiques ?}
  Les modèles traditionnels (\textit{ARIMA, Régression OLS}) ne sont pas adaptés aux infrastructures instables pour trois raisons :
  
  \vspace{0.5em}
  \begin{columns}[T]
    \begin{column}{0.48\textwidth}
      \begin{alertblock}{Surcharge de calcul}
        Calculer en continu sur une fenêtre glissante demande de stocker tout l'historique en mémoire vive et d'inverser de grandes matrices. C'est trop lourd pour des PC modestes.
      \end{alertblock}
    \end{column}
    
    \begin{column}{0.48\textwidth}
      \begin{alertblock}{Ignorance du contexte physique}
        Les pannes d'électricité (délestages) et les niveaux de batterie sont des signaux externes cruciaux. Les modèles classiques de séries temporelles ne savent pas les intégrer naturellement.
      \end{alertblock}
    \end{column}
  \end{columns}
\end{frame}

% Slide 2: Le Concept ARX
\begin{frame}{2. Qu'est-ce que le modèle ARX ?}
  Le modèle **ARX** combine deux sources d'informations complémentaires pour anticiper la disponibilité d'une machine :
  
  \vspace{0.5em}
  \begin{columns}[T]
    \begin{column}{0.48\textwidth}
      \begin{block}{La composante autorégressive (AR)}
        \begin{itemize}
          \item Utilise le passé récent de la machine (charges CPU et RAM de la minute précédente) pour prédire l'état futur immédiat.
        \end{itemize}
      \end{block}
    \end{column}
    
    \begin{column}{0.48\textwidth}
      \begin{block}{La composante exogène (X)}
        Intègre des signaux physiques directs de l'environnement :
        \begin{itemize}
          \item \textbf{Alimentation :} Statut secteur (branché/délestage) et batterie.
          \item \textbf{Réseau :} Connexion active ou coupée.
          \item \textbf{Humain :} Temps d'inactivité de l'utilisateur.
        \end{itemize}
      \end{block}
    \end{column}
  \end{columns}
\end{frame}

% Slide 3: L'Algorithme RLS
\begin{frame}{3. L'Algorithme RLS : L'apprentissage sans stockage}
  Le modèle **RLS** (\textit{Recursive Least Squares}) est le moteur d'apprentissage continu de notre système. Il met à jour le modèle à chaque minute sans stocker l'historique :
  
  \vspace{0.5em}
  \begin{columns}[T]
    \begin{column}{0.48\textwidth}
      \begin{block}{Fonctionnement en 3 étapes simples}
        \begin{enumerate}
          \item \textbf{Mesure de l'erreur :} Comparaison immédiate entre la prédiction et l'état réel de la machine.
          \item \textbf{Calcul du gain :} Dosage de la mise à jour (faut-il ajuster beaucoup ou peu les poids du modèle ?).
          \item \textbf{Mise à jour des poids :} Ajustement instantané de la mémoire interne.
        \end{enumerate}
      \end{block}
    \end{column}
    
    \begin{column}{0.48\textwidth}
      \begin{exampleblock}{Le Facteur d'Oubli ($\lambda$)}
        Un paramètre mathématique simple qui permet au modèle d'oublier progressivement les données très anciennes pour s'adapter rapidement aux nouveaux profils de délestages locaux.
      \end{exampleblock}
    \end{column}
  \end{columns}
\end{frame}

% Slide 4: Équivalence Temporelle de 72h
\begin{frame}{4. Comment régler la mémoire du modèle à 72 heures ?}
  Le professeur demande à juste titre une mise à jour sur une fenêtre glissante de 72 heures. Le RLS y parvient de manière élégante :
  
  \vspace{0.5em}
  \begin{columns}[T]
    \begin{column}{0.48\textwidth}
      \begin{block}{L'Équivalence Mathématique}
        \begin{itemize}
          \item Notre agent prend 1 mesure par minute.
          \item Une fenêtre de 72 heures représente exactement **4 320 mesures**.
          \item En ajustant le facteur d'oubli $\lambda$ à **0,99977**, le RLS accorde la même importance aux données qu'un historique glissant physique de 72h.
        \end{itemize}
      \end{block}
    \end{column}
    
    \begin{column}{0.48\textwidth}
      \begin{exampleblock}{Pourquoi c'est supérieur ?}
        \begin{itemize}
          \item \textbf{Économie de RAM :} Plus besoin de conserver les 4 320 points en mémoire.
          \item \textbf{Zéro interpolation :} Gère naturellement les déconnexions (si la machine s'éteint, l'apprentissage se fige sans planter).
        \end{itemize}
      \end{exampleblock}
    \end{column}
  \end{columns}
\end{frame}

% Slide 5: Analyse Comparative
\begin{frame}{5. OLS Glissant vs RLS Adaptatif}
  \begin{table}[]
    \centering
    \footnotesize
    \begin{tabular}{lll}
      \toprule
      \textbf{Critère de comparaison} & \textbf{Fenêtre Glissante (Méthode OLS)} & \textbf{Modèle RLS Adaptatif} \\
      \midrule
      \textbf{Usage de la RAM} & Conserve 4320 points de mesures & Conserve uniquement 2 matrices ($< 1$ Ko) \\
      \textbf{Vitesse d'exécution} & Ralentit à chaque mise à jour & Stable et instantané ($\le 0,1$ ms) \\
      \textbf{Calcul Matriciel} & Inversion lourde et risquée & Évitée (Algorithme stable par conception) \\
      \textbf{Gestion des Pannes} & Trous de données complexes à gérer & Gel automatique de l'apprentissage \\
      \textbf{Priorité Temporelle} & Traite le passé et le présent de façon égale & Donne naturellement la priorité au présent \\
      \bottomrule
    \end{tabular}
  \end{table}
\end{frame}

% Slide 6: Pourquoi le RLS est robuste aux anomalies ?
\begin{frame}{6. Immunité contre l'Oubli Catastrophique}
  Les réseaux de neurones peuvent souffrir d'un oubli brutal de leurs connaissances. Le modèle RLS-ARX résout ce problème grâce à deux verrous théoriques :
  
  \vspace{0.5em}
  \begin{columns}[T]
    \begin{column}{0.48\textwidth}
      \begin{block}{1. Convexité mathématique}
        \begin{itemize}
          \item Contrairement aux modèles profonds, le RLS possède une unique solution optimale globale. 
          \item Le modèle ne peut pas "sauter" d'un état stable à un état erratique lors d'une perturbation.
        \end{itemize}
      \end{block}
    \end{column}
    
    \begin{column}{0.48\textwidth}
      \begin{block}{2. Mémoire de la Covariance}
        \begin{itemize}
          \item Le modèle mémorise le niveau de certitude de ses prédictions.
          \item Face à un événement inhabituel ou transitoire (ex. micro-coupure réseau), il filtre automatiquement la perturbation sans altérer ses poids stables.
        \end{itemize}
      \end{block}
    \end{column}
  \end{columns}
\end{frame}

% Slide 7: Architecture Hybride
\begin{frame}{7. L'Architecture Prédictive Hybride de l'Agent}
  Pour une efficacité maximale, notre agent utilise une architecture à deux niveaux :
  
  \vspace{0.5em}
  \begin{columns}[T]
    \begin{column}{0.48\textwidth}
      \begin{block}{Le Réflexe Local Frugal : RLS-ARX}
        \begin{itemize}
          \item S'exécute à chaque minute.
          \item Réagit instantanément aux coupures électriques et variations de batterie locales.
          \item Consomme moins d'1 Ko de RAM.
        \end{itemize}
      \end{block}
    \end{column}
    
    \begin{column}{0.48\textwidth}
      \begin{block}{Le Cerveau Temporel : GRU + EWC}
        \begin{itemize}
          \item S'exécute périodiquement (ex. toutes les 24h).
          \item Apprend les habitudes d'usage de l'utilisateur sur le long terme.
          \item Protégé par l'algorithme EWC pour ne pas oublier ses connaissances globales.
        \end{itemize}
      \end{block}
    \end{column}
  \end{columns}
\end{frame}

% Slide 8: Frugalité Matérielle et TinyML
\begin{frame}{8. Intégration et Frugalité Matérielle (TinyML)}
  L'application des concepts TinyML sur nos nœuds volontaires à Yaoundé garantit le respect des ressources matérielles des participants :
  
  \vspace{0.5em}
  \begin{columns}[T]
    \begin{column}{0.48\textwidth}
      \begin{block}{Besoins en Mémoire RAM}
        \begin{itemize}
          \item Modèle linéaire RLS-ARX : **$< 1$ Ko**.
          \item Modèle GRU quantifié en INT8 : **$< 2$ Mo**.
          \item Totalement compatible avec la limite stricte de **10 Mo** pour les environnements contraints.
        \end{itemize}
      \end{block}
    \end{column}
    
    \begin{column}{0.48\textwidth}
      \begin{block}{Impact sur le Processeur (CPU)}
        \begin{itemize}
          \item Temps de calcul RLS : **$< 0,1$ milliseconde** par mise à jour.
          \item Inférence légère qui ne ralentit pas les tâches courantes de l'ordinateur de l'étudiant.
        \end{itemize}
      \end{block}
    \end{column}
  \end{columns}
\end{frame}

% Slide 9: Conclusion et Perspectives
\begin{frame}{9. Conclusion}
  \begin{block}{Les atouts majeurs du modèle RLS-ARX}
    Le choix du RLS-ARX répond parfaitement aux exigences académiques et pratiques :
    \begin{itemize}
      \item \textbf{Frugalité prouvée :} Convient parfaitement au TinyML sur des machines volontaires d'entrée de gamme.
      \item \textbf{Rigueur scientifique :} Équivalent exact à une fenêtre glissante de 72h grâce au facteur d'oubli calculé ($\lambda = 0,99977$).
      \item \textbf{Haute résilience :} Stable face à l'instabilité électrique et aux arrêts brusques inhérents au contexte de Yaoundé.
    \end{itemize}
  \end{block}
  
  \centering
  \vspace{1.5em}
  \Large\textbf{Merci pour votre attention. Place aux questions.}
\end{frame}

\end{document}
